\documentclass[11pt]{article}
\usepackage[margin=1in]{geometry}
\usepackage{amsmath,amssymb,amsfonts,amsthm,mathtools}
\usepackage{bm}
\usepackage{enumitem}
\usepackage[hidelinks]{hyperref}
\usepackage{cite}
\usepackage{mathtools}
\usepackage[numbers]{natbib} 
\usepackage[nameinlink,capitalise]{cleveref}
\usepackage{slashed}

% ===== Encoding & fonts (prevents Unicode errors) =====
\usepackage[T1]{fontenc}
\usepackage[utf8]{inputenc}
\usepackage{lmodern}
\usepackage{microtype}
\usepackage{authblk}


\newcommand{\Tr}{\mathrm{Tr}}
\newcommand{\Nil}{\mathrm{Nil}}
\newcommand{\Trgrav}{\mathrm{Tr}_{\mathrm{grav}}}
\newcommand{\Rplus}{R_{+}}
\newcommand{\ii}{\mathrm{i}}
\newcommand{\Kahler}{K\"ahler}
\newcommand{\J}{J}
\newcommand{\Om}{\Omega}
\newcommand{\LensNil}{\text{Lens}\times\text{Nil}}
\newcommand{\dd}{\mathrm{d}}
\newcommand{\bbZ}{\mathbb{Z}}
\newcommand{\bbC}{\mathbb{C}}
\newcommand{\cO}{\mathcal{O}}
\newcommand{\vol}{\mathrm{vol}}
\newcommand{\MTT}{Modal Triplet Theory}
\newcommand{\Hcoh}{\mathcal{H}_{\text{coh}}}
\newcommand{\Dprop}{\Delta_{\mathrm{prop}}}
\newcommand{\DBV}{\Delta_{\mathrm{BV}}}
\newcommand{\Fform}{F}
\newcommand{\LSPT}{\Lambda}
\newcommand{\lamgap}{\lambda_{\ast}}
\newcommand{\alphap}{\alpha'}
\newcommand{\RR}{\mathbb{R}}
\newcommand{\tr}{\mathrm{tr}}
\newcommand{\Vol}{\mathrm{Vol}}
\newcommand{\astone}{\ast\!1}
\newcommand{\R}{\mathbb{R}}
\newcommand{\C}{\mathbb{C}}
\newcommand{\cH}{\mathcal{H}}
\newcommand{\Id}{\mathbbm{1}}
\newcommand{\wh}[1]{\widehat{#1}}
\newcommand{\Hobs}{\mathcal{H}_{\mathrm{obs}}}
\newcommand{\HQM}{\mathcal{H}_{\mathrm{QM}}}
\newcommand{\Leff}{\mathcal{L}_{\mathrm{eff}}}
\newcommand{\CC}{\mathbb{C}}
\newcommand{\TT}{\mathbb{T}}
\newcommand{\Dom}{\mathrm{Dom}}
\newcommand{\Ran}{\mathrm{Ran}}


\DeclareMathOperator{\Ric}{Ric}

\let\varphi\phi

\newtheorem{theorem}{Theorem}[section]
\newtheorem{lemma}[theorem]{Lemma}
\newtheorem{definition}[theorem]{Definition}
\newtheorem{proposition}[theorem]{Proposition}
\newtheorem{corollary}[theorem]{Corollary}
\newtheorem{assumption}[theorem]{Assumption}
\newtheorem{conjecture}[theorem]{Conjecture}
\newtheorem{example}[theorem]{Example}
\newtheorem{remark}[theorem]{Remark}
\newtheorem{note}[theorem]{Note}

% --- Custom Commands ---

\theoremstyle{definition}
\theoremstyle{plain}
\theoremstyle{remark}

% cross-referencing helpers
\newcommand{\secref}[1]{Section~\ref{#1}}
\newcommand{\appref}[1]{Appendix~\ref{#1}}
\newcommand{\thmref}[1]{Theorem~\ref{#1}}
\newcommand{\propref}[1]{Proposition~\ref{#1}}
\newcommand{\lemref}[1]{Lemma~\ref{#1}}
\newcommand{\ksym}{\texorpdfstring{$\kappa$}{kappa}\nobreakdash-symmetry}

\numberwithin{equation}{section}

\title{\textbf{Modal Triplet Theory: Quantum Amplitudes from Modal Geometry}\\
\large A First--Principles Derivation of Feynman Rules, Running, and Phenomenology}
\author{Peter Nero}
\date{September 8, 2025}

\begin{document}
\maketitle

\begin{abstract}
We present a first--principles derivation of perturbative quantum field theory amplitudes
from Modal Triplet Theory (MTT). Starting from the coherent fixed--point sector of the
modal configuration space, we employ the MTT $\rightarrow$ QFT projection to obtain a
locally covariant algebraic quantum field theory (AQFT) on a globally hyperbolic
four--dimensional spacetime. Within this framework we construct the local net of
observables, establish algebraic propagation via the time--slice axiom, and define
interacting quantum fields using perturbative algebraic quantum field theory (pAQFT).

Feynman rules, loop corrections, and renormalization are derived as consequences of the
algebraic construction rather than postulated. Dimensional regularization with minimal
subtraction and the covariant functional renormalization group are recovered as concrete
realizations of the underlying algebraic renormalization freedom. We identify the precise
conditions under which scattering amplitudes exist, formulate amplitudes via
Haag--Ruelle/LSZ theory in admissible scattering regimes, and distinguish them from local
or in--in observables in generic curved or time--dependent backgrounds.

All effective couplings, masses, and mixing parameters are determined by bounded overlap
integrals and curvature--gap data intrinsic to Modal Triplet Theory. When evaluated with
identical low--energy inputs, the resulting amplitudes coincide with those of the Standard
Model. The construction yields a complete and logically closed amplitudes layer for MTT,
bridging modal geometry to perturbative quantum field theory in a mathematically rigorous
and regime--aware manner.
\end{abstract}


\tableofcontents

\section{Introduction and scope}

Modal Triplet Theory (MTT) provides a fixed--point, gap--structured framework in which
four--dimensional physics emerges from a higher--dimensional modal geometry. Previous
work has established explicit projections from MTT to quantum mechanics and to general
relativity, as well as a projection to algebraic quantum field theory (QFT) on curved
spacetimes. In that construction, the projection $\Pi_{\mathrm{QFT}}$ maps a coherent
modal configuration to a locally covariant algebraic QFT equipped with a Hadamard state,
thereby enabling well--defined local observables, renormalization, and backreaction.

The purpose of the present paper is to complete the program by deriving
\emph{perturbative quantum field theory amplitudes} from Modal Triplet Theory in a manner
that is mathematically rigorous, logically closed, and explicit about its domain of
validity. In particular, we address the following question:

\begin{quote}
Given a coherent modal sector and the induced algebraic QFT on spacetime, under what
conditions do standard Feynman rules and scattering amplitudes emerge, and how are they
computed from first principles?
\end{quote}

A central theme of this work is that scattering amplitudes are \emph{not} fundamental
objects. In generic curved or time--dependent spacetimes, no global $S$--matrix exists,
and physical predictions must be formulated in terms of local or in--in observables.
Amplitudes arise only in special regimes admitting asymptotic particle states, such as
asymptotically flat or stationary backgrounds. Any first--principles derivation of
amplitudes must therefore make these regimes explicit rather than assuming them implicitly.

To achieve this, we proceed in four logically distinct steps.

First, we specify the Modal Triplet Theory input data relevant for quantum field theory:
the coherent sector, the emergent four--dimensional field content, and the overlap
integrals that determine effective couplings, masses, and mixing parameters. No quantum
field theoretic assumptions enter at this stage.

Second, we formulate the kinematics of quantum fields using algebraic quantum field theory.
We construct the local net of observables, establish microcausality, and identify the
time--slice axiom as the precise algebraic statement of causal propagation. This step
provides the correct notion of locality and dynamics without reference to particles or
Feynman diagrams.

Third, we introduce interactions using perturbative algebraic quantum field theory
(pAQFT). Interacting fields are defined via time--ordered products and the Bogoliubov
map, with renormalization appearing as a finite and local ambiguity classified by
covariance and power counting. Standard schemes such as dimensional regularization with
minimal subtraction and the covariant functional renormalization group are recovered as
specific realizations of this algebraic freedom.

Fourth, we show how Feynman rules and scattering amplitudes emerge from the algebraic
framework. Wick expansion of time--ordered products yields the familiar diagrammatic
representation, while Haag--Ruelle and LSZ theory identify the precise conditions under
which amputated correlation functions define scattering amplitudes. Outside these
conditions, the same formalism yields local correlation functions or in--in observables.

The result is a complete amplitudes layer for Modal Triplet Theory. All perturbative
quantum field theoretic structures---propagators, vertices, loop corrections, running
couplings, and benchmark amplitudes---are derived rather than assumed, and their regime
of validity is stated explicitly. When evaluated with the same low--energy inputs, the
resulting amplitudes coincide with those of the Standard Model, while retaining a clear
geometric origin in the modal overlaps and curvature--gap data.

The paper is organized as follows. In Section~2 we summarize the Modal Triplet Theory input
data and the QFT projection. Section~3 derives the effective four--dimensional field
content and the overlaps$\rightarrow$parameters pipeline. Sections~4--7 develop the
algebraic and perturbative quantum field theory framework, culminating in the derivation
of Feynman rules. Section~8 defines scattering amplitudes and their regimes of validity.
Section~9 presents benchmark calculations. Subsequent sections address gauge symmetry,
anomalies, confinement bounds, and the synthesis of this work with the broader MTT corpus.

\section{Modal Triplet Theory input and the QFT projection}

This section summarizes the Modal Triplet Theory (MTT) structures that serve as input for
the quantum field theoretic construction developed in the remainder of the paper. No new
results are proved here; rather, we isolate the precise data inherited from MTT and fix
notation. Readers are referred to the foundational and projection papers in the MTT corpus
for detailed proofs.

\subsection{Modal configuration space and coherent sector}

The kinematical configuration space of Modal Triplet Theory is a ten--dimensional manifold
of the form
\[
M_{10} \;=\; Y_4 \times B_1 \times B_2 \times B_3 ,
\]
where $Y_4$ is a four--dimensional Lorentzian spacetime and the $B_n$ are compact internal
modal bundles. Each $B_n$ carries its own elliptic operator (typically a Laplacian or
Dirac--type operator) governing internal excitations.

MTT dynamics admit stable coherent fixed points characterized by a spectral gap separating
a finite--dimensional coherent sector from higher internal excitations. Let
$\Pi_{\mathrm{coh}}$ denote the joint Riesz projector onto this coherent sector. The
existence, boundedness, and stability of $\Pi_{\mathrm{coh}}$ on admissible spacetime slabs
are guaranteed by the Fundamental Contractivity Condition and the fixed--point theorems of
the MTT foundation.

All constructions in this paper are conditional on working within such an admissible slab,
where coherent projection is well defined and stable. No global--in--time extension is
assumed.

\subsection{Effective four--dimensional field content}

Restriction of a coherent modal configuration to the base spacetime $Y_4$ yields an
effective four--dimensional field content consisting of:
\begin{itemize}
\item gauge fields associated with internal isometries or symmetry bundles,
\item fermionic fields arising from spinor representations on the internal bundles,
\item scalar fields corresponding to coherent internal modes and symmetry breaking
sectors.
\end{itemize}

The representation content, family structure, and anomaly cancellation properties are
determined by internal topology, holonomy, and the ``two--of--three'' selection mechanism
characteristic of MTT. These features fix the field multiplets and symmetry assignments of
the effective theory independently of any quantum field theoretic considerations.

\subsection{Overlap integrals and geometric parameters}

All dimensionless couplings and mass parameters of the effective four--dimensional theory
are determined by bounded overlap integrals on the internal bundles. In particular:
\begin{itemize}
\item gauge couplings arise from quadratic overlaps of normalized internal harmonic forms,
\item Yukawa couplings arise from trilinear overlaps,
\item mixing angles and CP--violating phases arise from internal holonomy data,
\item effective mass scales and thresholds receive curvature--dependent shifts governed by
representation--correct gap laws.
\end{itemize}

These overlap integrals are finite by construction and depend only on the modal geometry.
They provide numerical input data for the quantum field theory but do not rely on any
perturbative or diagrammatic machinery.

\subsection{The QFT projection \texorpdfstring{$\Pi_{\mathrm{QFT}}$}{PiQFT}}

The projection from Modal Triplet Theory to quantum field theory is implemented by a map
\[
\Pi_{\mathrm{QFT}} :
\{\text{coherent modal configurations}\}
\;\longrightarrow\;
(\mathcal{A}(Y_4,g),\omega) ,
\]
where $\mathcal{A}(Y_4,g)$ is an algebra of quantum fields on the spacetime $(Y_4,g)$ and
$\omega$ is a state on that algebra.

The defining properties of $\Pi_{\mathrm{QFT}}$ are:
\begin{itemize}
\item \textbf{Algebraic output:} The image is a CCR (bosonic) and CAR (fermionic) algebra of
quantum fields, not a Hilbert space representation.
\item \textbf{Hadamard state selection:} The state $\omega$ satisfies the Hadamard condition,
ensuring that local composite operators and renormalized expectation values are well
defined.
\item \textbf{Local covariance:} The construction is covariant under admissible embeddings
of spacetime regions.
\item \textbf{Slab--local validity:} The projection is defined only on finite spacetime slabs
where coherent projection and bounded geometry hold.
\end{itemize}

The QFT projection incorporates no additional dynamical assumptions beyond those already
present in MTT. In particular, it does not assume a Lagrangian path integral, Feynman
diagrams, or scattering theory. It supplies only the algebraic and state--theoretic data
required for quantum field theory.

\subsection{Role of this section}

The data specified above constitute the \emph{entire input} from Modal Triplet Theory used
in the remainder of the paper. All subsequent constructions---local nets, interacting
fields, renormalization, and amplitudes---are derived from these inputs using algebraic
quantum field theory and its perturbative extension.

In particular, no phenomenological assumptions or standard quantum field theoretic rules
are imported at this stage. The emergence of Feynman rules and amplitudes will be shown
explicitly in later sections.
\section{Emergent $3{+}1$ field content and the overlaps$\rightarrow$parameters pipeline}

In this section we derive the effective four--dimensional field theory that serves as input
to the algebraic quantum field theoretic construction developed in later sections. The
result is a canonically normalized classical action on $(Y_4,g)$ together with a complete
set of parameters determined by modal geometry. No quantum field theoretic assumptions
enter at this stage.

\subsection{Emergent classical action}

Restriction of a coherent modal configuration to the base spacetime $Y_4$ yields an
effective classical action of the form
\[
S_{\mathrm{eff}}
=
\int_{Y_4} d^4x\,\sqrt{-g}\,
\Big[
-\frac{1}{4} F^a_{\mu\nu} F^{a\,\mu\nu}
+ \bar{\psi}\, i \gamma^\mu D_\mu \psi
+ |D_\mu \phi|^2
- V(\phi)
- \bar{\psi}_i Y_{ij} \psi_j \phi
+ \mathrm{h.c.}
\Big],
\]
where $F^a_{\mu\nu}$ are gauge field strengths, $\psi$ denotes fermionic multiplets,
$\phi$ scalar fields, and $D_\mu$ the gauge--covariant derivative determined by the
representation content fixed in Section~2.

The field content and symmetry assignments are inherited directly from the internal modal
geometry. In particular, the gauge group, matter representations, and family structure are
determined independently of any quantum field theoretic considerations.

\subsection{Canonical normalization}

The effective action obtained by dimensional reduction is canonically normalized by
construction. Internal harmonic representatives on the modal bundles $B_n$ are chosen to
be orthonormal with respect to the natural $L^2$ inner product induced by the internal
metric. With this choice:
\begin{itemize}
\item gauge kinetic terms reduce to the standard Yang--Mills form,
\item fermionic kinetic terms reduce to the standard Dirac form,
\item scalar kinetic terms are canonically normalized.
\end{itemize}

All dimensionful factors arising from the higher--dimensional action and internal volumes
are absorbed into the normalization of the effective couplings defined below. No further
field rescalings are required at the four--dimensional level.

\subsection{Gauge couplings from overlap integrals}

Gauge couplings are determined by quadratic overlap integrals of normalized internal modes.
For each gauge factor labelled by $a$, one has
\[
\frac{1}{g_a^2}
=
N_a \int_{B_a} \langle \omega_a, \omega_a \rangle \, d\mu_{B_a},
\]
where $\omega_a$ is the normalized harmonic representative on the corresponding internal
bundle and $N_a$ is a fixed numerical factor determined by the higher--dimensional action
and representation conventions.

For hypercharge we adopt the GUT--normalized convention
\[
g_1 \equiv \sqrt{\frac{5}{3}}\, g',
\qquad
\alpha_1 = \frac{5}{3}\,\alpha_Y,
\]
so that the one--loop $\beta$--function coefficient takes its standard value. All gauge
couplings are finite and computable from the modal geometry.

\subsection{Yukawa couplings, mixing, and CP phases}

Yukawa couplings arise from trilinear overlap integrals on the full internal space
$B_1 \times B_2 \times B_3$. Their magnitudes are bounded by construction, while flavor
mixing angles and CP--violating phases are determined by internal holonomy data, in
particular by phases associated with nontrivial cycles of the internal geometry.

The resulting Yukawa matrices are generic subject to the symmetry and selection rules of
Modal Triplet Theory, and provide the complete set of fermion masses and mixings after
spontaneous symmetry breaking.

\subsection{Curvature--gap mass shifts and thresholds}

In addition to masses generated by spontaneous symmetry breaking, effective mass scales
receive curvature--dependent shifts governed by representation--correct gap laws. For
example, in a background with scalar curvature $\mathrm{Scal}$,
\begin{align}
M^2_{\mathrm{spinor}} &\sim m^2 + \frac{1}{4}\,\mathrm{Scal}, \\
M^2_{\mathrm{scalar}} &\sim m^2 + \frac{1}{6}\,\mathrm{Scal},
\end{align}
with analogous expressions for higher--spin fields.

These curvature--gap masses determine physical thresholds in the effective field theory
and play a central role in decoupling and renormalization group flow. All such thresholds
are fixed geometrically and introduce no free parameters.

\subsection{Summary}

At the end of this section, we have obtained:
\begin{itemize}
\item a canonically normalized classical field theory on $(Y_4,g)$,
\item a complete set of gauge, Yukawa, and scalar couplings determined by overlap integrals,
\item geometrically fixed mass scales and thresholds.
\end{itemize}

No quantum field theoretic constructions have yet been introduced. In the next section we
use this classical data as input to algebraic quantum field theory, where locality and
propagation are defined independently of perturbative diagrams or scattering theory.


\section{Algebraic Quantum Field Theory: Kinematics}

In this section we define the kinematic framework within which quantum fields propagate.
Propagation is formulated algebraically, without reference to particle trajectories or
background transport laws. This provides the precise notion of locality and causality
required for a first--principles derivation of amplitudes.

Throughout this section, let $(Y_4,g)$ denote the globally hyperbolic four--dimensional
spacetime obtained from the Modal Triplet Theory projection $\Pi_{\mathrm{QFT}}$ on a
fixed admissible slab. All statements are understood slab--locally; no global--in--time
claims are made.

\subsection{Field algebras on curved spacetime}

Let $P$ denote a normally hyperbolic operator defining the free equation of motion for a
bosonic field $\phi$ (e.g.\ Klein--Gordon with curvature coupling), or a Dirac--type operator
for fermionic fields. Let $\Delta = E_{\mathrm{ret}} - E_{\mathrm{adv}}$ be the causal
propagator associated with $P$.

The algebra of observables $\mathcal{A}(Y_4,g)$ is defined as the CCR (bosonic) or CAR
(fermionic) $*$--algebra generated by smeared fields $\Phi(f)$ (or $\Psi(f)$) with
$f \in C^\infty_0(Y_4)$, modulo:
\begin{align}
[\Phi(f),\Phi(h)] &= i\Delta(f,h)\,\mathbf{1}, \\
\Phi(Pf) &= 0 ,
\end{align}
and the analogous anticommutation relations for fermions.

States on $\mathcal{A}(Y_4,g)$ are positive normalized linear functionals. In the MTT
projection, the state $\omega$ selected by $\Pi_{\mathrm{QFT}}$ is of Hadamard form,
ensuring well--defined local observables and renormalization.

\subsection{Local nets}

For each open region $\mathcal{O}\subset Y_4$, define the local algebra
\[
\mathcal{A}(\mathcal{O}) :=
\langle \Phi(f)\;|\; f\in C^\infty_0(Y_4),\ \mathrm{supp}(f)\subset\mathcal{O} \rangle_* .
\]

This assignment satisfies the Haag--Kastler axioms relevant for the present work:

\paragraph{Isotony.}
If $\mathcal{O}_1 \subset \mathcal{O}_2$, then
$\mathcal{A}(\mathcal{O}_1) \subset \mathcal{A}(\mathcal{O}_2)$.

\paragraph{Microcausality.}
If $\mathcal{O}_1$ and $\mathcal{O}_2$ are causally disjoint, then
\[
[\mathcal{A}(\mathcal{O}_1),\mathcal{A}(\mathcal{O}_2)] = 0
\]
for bosons (or graded commutator zero for fermions).

These properties follow directly from the support properties of the causal propagator
$\Delta$.

\subsection{Time--slice axiom and algebraic propagation}

The fundamental propagation principle in algebraic quantum field theory is the time--slice
axiom.

\begin{theorem}[Time--slice property]
Let $\mathcal{O}\subset Y_4$ be an open region containing a Cauchy surface for $(Y_4,g)$
(or for the relevant admissible slab). Then the inclusion
\[
\mathcal{A}(\mathcal{O}) \hookrightarrow \mathcal{A}(Y_4,g)
\]
is an isomorphism onto the algebra of the domain of dependence $D(\mathcal{O})$.
\end{theorem}

This theorem provides the precise algebraic meaning of causal propagation: all observables
in the causal future and past of $\mathcal{O}$ are generated by data localized arbitrarily
close to a Cauchy surface. No additional degrees of freedom propagate superluminally or
nonlocally.

The AQFT net $(\mathcal{O}\mapsto\mathcal{A}(\mathcal{O}))$ constructed here coincides with
the net induced by admissible chart overlap in Modal Triplet Theory when the QFT projection
exists, as established in the companion work on AQFT nets from MTT.
\section{Interacting Fields via Perturbative Algebraic QFT}

We now construct interacting quantum fields without assuming a Lagrangian path integral
or diagrammatic rules. Interactions are defined algebraically using perturbative algebraic
quantum field theory (pAQFT), which is compatible with curved spacetimes and local covariance.

\subsection{Local functionals and observables}

Let $\mathcal{F}_{\mathrm{loc}}$ denote the space of local classical functionals of the field
configuration, i.e.\ functionals of the form
\[
F(\phi) = \int_{Y_4} \mathcal{L}(x,\phi(x),\nabla\phi(x),\ldots)\, d\mathrm{vol}_g ,
\]
with compact support and polynomial dependence on the fields and their derivatives.

Quantum observables are obtained by deformation quantization of $\mathcal{F}_{\mathrm{loc}}$
using a chosen Hadamard two--point function, inducing a noncommutative $\star$--product.

\subsection{Time--ordered products}

A renormalized time--ordered product is a family of multilinear maps
\[
T_n : \mathcal{F}_{\mathrm{loc}}^{\otimes n}
\longrightarrow \mathcal{A}(Y_4,g)[[\hbar]]
\]
satisfying:
\begin{itemize}
\item locality and covariance,
\item symmetry under permutations,
\item causal factorization,
\item the microlocal spectrum condition.
\end{itemize}

For Hadamard states on globally hyperbolic spacetimes, such time--ordered products exist,
with a finite renormalization freedom classified by local covariant counterterms
(Hollands--Wald).

\subsection{Interacting algebra}

Let $V \in \mathcal{F}_{\mathrm{loc}}$ denote the interaction functional corresponding to the
effective interactions determined by Modal Triplet Theory (gauge couplings, Yukawas,
scalar potentials).

Define the relative $S$--matrix
\[
S(V) := T \exp\!\left(\frac{i}{\hbar} V\right)
= \sum_{n\ge0} \frac{i^n}{n!\hbar^n} T_n(V^{\otimes n}) .
\]

Interacting observables are defined via the Bogoliubov map
\[
R_V(F) :=
\left.\frac{\hbar}{i}\frac{d}{d\lambda}\right|_{\lambda=0}
S(V)^{-1} \star S(V+\lambda F) .
\]

The interacting algebra $\mathcal{A}_V(\mathcal{O})$ is generated by $R_V(F)$ with
$\mathrm{supp}(F)\subset\mathcal{O}$. Causal factorization guarantees that
$\mathcal{O}\mapsto\mathcal{A}_V(\mathcal{O})$ again defines a local net satisfying isotony
and microcausality.

This construction defines interacting quantum fields in a mathematically controlled way,
independent of perturbative diagrams.
\section{Renormalization from First Principles}

Renormalization in the present framework is not introduced as a procedural subtraction,
but arises as a structural feature of the algebraic definition of time--ordered products.

\subsection{Renormalization freedom}

The ambiguity in the definition of the time--ordered products $T_n$ is finite and local,
corresponding to the addition of local covariant counterterms consistent with:
\begin{itemize}
\item locality and covariance,
\item power counting,
\item symmetries (gauge, BRST, discrete symmetries).
\end{itemize}

This freedom is precisely the renormalization freedom classified in curved--spacetime
QFT (Hollands--Wald).

\subsection{Dimensional regularization and minimal subtraction}

In flat spacetime or approximately flat regimes, dimensional regularization with
minimal subtraction provides a concrete realization of the algebraic renormalization
freedom. The counterterms selected by $\overline{\mathrm{MS}}$ correspond to a particular
choice of $T_n$ within the allowed class.

The one--loop $\beta$--functions and anomalous dimensions used in later sections are
instances of this general algebraic structure.

\subsection{Covariant functional renormalization group}

To track scale dependence on curved spacetimes, we employ the covariant functional
renormalization group (FRG). The effective average action $\Gamma_k$ satisfies the
Wetterich equation with a background--covariant regulator.

In the MTT framework, heavy internal modes generate threshold corrections determined
by curvature--gap masses. As long as the admissibility margins remain positive, the FRG
flow preserves locality and the validity of the effective field theory.

Thus DR/MS and FRG are understood as complementary realizations of the same underlying
renormalization freedom inherent in the algebraic construction.
\section{Emergence of Feynman Rules}

We now show how the familiar Feynman rules arise from the algebraic construction of the
preceding sections.

\subsection{Wick expansion in Hadamard states}

For quasifree Hadamard states, Wick's theorem holds for the free theory. Products of
fields can be expanded in terms of normal--ordered products and contractions determined
by the Hadamard two--point function.

\subsection{Graph expansion of time--ordered products}

Expanding the relative $S$--matrix and the Bogoliubov map in powers of the interaction
functional $V$ yields a formal power series whose coefficients are finite sums over
graphs:
\begin{itemize}
\item vertices correspond to local monomials in $V$,
\item propagators correspond to the two--point function,
\item symmetry factors arise from combinatorics of $T_n$.
\end{itemize}

Renormalization amounts to replacing ill--defined graph contributions by the
renormalized time--ordered products selected in Sec.~6.

\subsection{Equivalence with standard perturbative rules}

In asymptotically flat regimes, the resulting expansion coincides with the standard
Feynman rules of perturbative quantum field theory. The propagators and vertices listed
in Sec.~3.3 and Appendix~A are therefore derived consequences of the algebraic framework,
not independent assumptions.

This establishes the logical bridge from Modal Triplet Theory to conventional
perturbative amplitudes.
\section{Scattering Theory and the Definition of Amplitudes}
\label{sec:scattering-amplitudes}

Having constructed the local net of observables, the interacting algebra, and the
perturbative expansion from first principles, we now address the precise meaning of
\emph{scattering amplitudes}. This requires distinguishing regimes in which asymptotic
particle states exist from those in which only local or in--in observables are well defined.

\subsection{Scattering regimes}

Scattering theory requires more structure than local propagation. In particular, one
needs:
\begin{itemize}
\item a notion of asymptotic time evolution,
\item a stable one--particle subspace,
\item a mass gap separating single--particle states from the continuum.
\end{itemize}

Accordingly, we restrict the notion of scattering amplitudes to regimes in which the
effective spacetime $(Y_4,g)$ is asymptotically stationary or asymptotically flat on the
admissible slab. This includes Minkowski space, stationary curved backgrounds, and
cosmological spacetimes with asymptotically adiabatic regions.

Outside such regimes, no global $S$--matrix exists, and amplitudes must be replaced by
local correlation functions or in--in observables.

\subsection{Asymptotic fields and Haag--Ruelle theory}

Assume that $(Y_4,g)$ admits asymptotic regions in which the interacting dynamics
approaches a free theory with the same mass spectrum. Let $\mathcal{H}$ be the Hilbert
space obtained from a Hadamard state $\omega$ via the GNS construction.

Under the usual spectral and stability assumptions (mass gap, absence of infraparticles),
Haag--Ruelle theory applies. One may define asymptotic in-- and out--fields
\[
\Phi_{\mathrm{in/out}}(f)
=
\lim_{t\to\mp\infty} \Phi_t(f_t),
\]
where $f_t$ are suitably chosen solutions of the free field equation with compact
momentum support.

These asymptotic fields generate Fock spaces
$\mathcal{H}_{\mathrm{in}}$ and $\mathcal{H}_{\mathrm{out}}$ describing incoming and
outgoing particle states. The interacting dynamics defines a unitary scattering operator
\[
S:\mathcal{H}_{\mathrm{in}}\to\mathcal{H}_{\mathrm{out}},
\]
intertwining the asymptotic representations.

\subsection{LSZ reduction}

In scattering regimes, matrix elements of the $S$--matrix can be expressed in terms of
time--ordered correlation functions via LSZ reduction. For scalar fields, one has
schematically
\[
\langle p_1',\ldots,p_m'|S|p_1,\ldots,p_n\rangle
=
\prod_i (p_i^2-m^2)
\prod_j (p_j'^2-m^2)
\,
\widetilde{G}_{m+n}(p',p),
\]
where $\widetilde{G}_{m+n}$ denotes the Fourier transform of the renormalized
time--ordered $n{+}m$--point function constructed in
Secs.~5--7.

Field--strength renormalization factors are absorbed into the LSZ residues, which in the
present framework are determined by the normalization of the two--point function selected
by the MTT projection.

Thus, when the LSZ assumptions hold, the perturbative expansion of the interacting algebra
yields well--defined scattering amplitudes in the usual sense.

\subsection{Amplitudes versus local observables}

It is important to emphasize that scattering amplitudes are \emph{not} fundamental
objects in the present framework. They are derived quantities, available only in special
regimes. The fundamental objects are:
\begin{itemize}
\item local algebras $\mathcal{A}(\mathcal{O})$,
\item interacting observables $R_V(F)$,
\item local correlation functions.
\end{itemize}

In generic curved or time--dependent spacetimes where asymptotic states do not exist,
physical predictions are formulated in terms of in--in expectation values
\[
\langle R_V(F_1)\cdots R_V(F_n)\rangle_\omega ,
\]
computed using the same perturbative machinery. Cosmological particle production,
nonequilibrium processes, and backreaction effects fall into this category.

\subsection{Summary}

In summary:
\begin{itemize}
\item The AQFT and pAQFT constructions provide a universal definition of interacting
quantum fields and local observables.
\item Scattering amplitudes emerge only in regimes admitting asymptotic particle states.
\item When such regimes exist, LSZ and Haag--Ruelle theory identify amplitudes with
appropriately amputated time--ordered correlation functions.
\item Outside scattering regimes, the same formalism yields local or in--in observables
rather than a global $S$--matrix.
\end{itemize}

With this distinction in place, we now proceed to explicit benchmark calculations in
regimes where scattering amplitudes are well defined.
\section{Benchmark Calculations}
\label{sec:benchmarks}

In this section we present representative benchmark calculations demonstrating how
standard perturbative results emerge from the algebraic framework developed in
Secs.~4--8. All results below are derived from the interacting algebra and its
perturbative expansion, under the scattering assumptions specified in
Sec.~\ref{sec:scattering-amplitudes}. No additional dynamical or kinematic assumptions
are introduced at this stage.

\subsection{Tree--level scattering processes}

We begin with simple $2\to2$ scattering processes in asymptotically flat regimes.
In such settings, the LSZ conditions of Sec.~\ref{sec:scattering-amplitudes} apply and
scattering amplitudes are well defined.

\paragraph{$e^+e^- \to \mu^+\mu^-$.}
At energies far below the $Z$ pole, the dominant contribution arises from $s$--channel
photon exchange. The amplitude is obtained from the amputated four--point function
constructed via the perturbative expansion of the interacting algebra, yielding
\[
\sigma(s) = \frac{4\pi\alpha^2}{3s},
\]
with $\alpha$ determined by the overlap--defined gauge coupling at the reference scale.
Near the $Z$ pole, $Z$ exchange and $\gamma$--$Z$ interference enter, with couplings
$g_1$ and $g_2$ fixed by internal overlap integrals as described in
Sec.~\ref{sec:overlaps-parameters}.

\paragraph{$e^-\mu^- \to e^-\mu^-$.}
$t$--channel photon exchange in QED yields the standard Rutherford scattering amplitude.
Field--strength renormalization factors are absorbed into LSZ residues and do not alter
the tree--level result.

These examples illustrate that once canonical normalization is fixed by the overlap
construction, the resulting amplitudes coincide with their Standard Model counterparts.

\subsection{Weak charged--current processes}

\paragraph{$u\bar d \to W^+ \to \mu^+\nu_\mu$.}
Charged--current amplitudes are generated by the interaction functional $V$ containing
the $SU(2)_L$ gauge coupling. The amplitude is controlled by $g_2$, determined by overlap
integrals on $B_2$, and by the CKM matrix elements arising from internal holonomy data.
At tree level, the resulting cross section agrees with the Standard Model expression
when evaluated with the same low--energy inputs.

\subsection{One--loop running and threshold matching}

Loop corrections arise from the renormalized time--ordered products constructed in
Sec.~6. The one--loop renormalization group equations take the form
\[
\mu\frac{d g_a}{d\mu} = \frac{b_a}{16\pi^2} g_a^3 + \cdots ,
\]
with coefficients $(b_1,b_2,b_3) = (41/10,-19/6,-7)$ for
$U(1)_Y$, $SU(2)_L$, and $SU(3)_c$, respectively, using GUT--normalized hypercharge.

Threshold matching is implemented at scales $\mu \simeq M_{\mathrm{heavy}}$, where
$M_{\mathrm{heavy}}$ are the effective masses determined by curvature--gap eigenvalues
and heavy internal excitations. This realizes Appelquist--Carazzone decoupling within
the algebraic renormalization framework.

\subsection{Leptonic and pseudoscalar decays}

\paragraph{Muon decay.}
The muon lifetime is obtained from the effective four--fermion interaction generated
by integrating out the $W$ boson. The decay rate
\[
\Gamma_\mu =
\frac{G_F^2 m_\mu^5}{192\pi^3}
f\!\left(\frac{m_e^2}{m_\mu^2}\right)
\left[1+\Delta_{\mathrm{RC}}^{(\mu)}\right]
\]
is reproduced, with the radiative correction $\Delta_{\mathrm{RC}}^{(\mu)}$ arising from
one--loop time--ordered products. The Fermi constant $G_F$ is fixed by the overlap--defined
$SU(2)_L$ coupling and the $W$ mass.

\paragraph{$\pi^\pm,K^\pm \to \ell^\pm\nu_\ell[\gamma]$.}
The decay widths
\[
\Gamma(P\to \ell\nu[\gamma]) =
\frac{G_F^2 |V_{qQ}|^2 f_P^2 m_\ell^2 m_P}{8\pi}
\left(1-\frac{m_\ell^2}{m_P^2}\right)^2
\left[1+\delta_{\mathrm{EM}}^P+\delta_{\mathrm{SU(2)}}^P\right]
\]
are obtained, including electromagnetic and isospin--breaking corrections. In the
present framework, decay constants and CKM elements may be taken as empirical inputs
or evaluated directly from overlap integrals, yielding agreement with Standard Model
phenomenology when the same low--energy data are used.

\subsection{Non--scattering regimes}

In cosmological or otherwise time--dependent backgrounds where asymptotic states do not
exist, the same perturbative machinery yields local or in--in observables rather than
scattering amplitudes. Particle production in expanding spacetimes, nonequilibrium
effects, and backreaction phenomena are computed as expectation values in the interacting
algebra, using the same renormalized time--ordered products.

\subsection{Summary}

The benchmark calculations presented in this section demonstrate that:
\begin{itemize}
\item the algebraic construction of interacting fields reproduces standard perturbative
results in scattering regimes,
\item loop corrections and running couplings arise naturally from algebraic
renormalization,
\item thresholds are fixed by geometric curvature--gap data,
\item phenomenological agreement with the Standard Model is obtained when evaluated
with identical low--energy inputs.
\end{itemize}

All results are consequences of the AQFT and pAQFT framework developed in
Secs.~4--8, together with the Modal Triplet Theory determination of effective field
content and parameters.

\section{Gauge symmetry, BRST, and unitarity}
\label{sec:brst-unitarity}

Gauge symmetry plays a dual role in the present framework. At the classical level it is
fixed by the modal geometry and internal symmetry data of Modal Triplet Theory. At the
quantum level it constrains the admissible renormalization freedom and guarantees
unitarity of the interacting theory. In this section we show how gauge fixing, BRST
symmetry, and unitarity are implemented consistently within the algebraic and perturbative
framework developed in the preceding sections.

\subsection{Gauge fixing in the algebraic framework}

Gauge theories require gauge fixing in order to define propagators and time--ordered
products. In the algebraic setting this is implemented by enlarging the field algebra to
include ghost fields and auxiliary fields, and by working with a gauge--fixed action that
preserves locality and covariance.

We employ background--covariant gauge fixing conditions of Lorenz type,
\[
\mathcal{L}_{\mathrm{gf}}
=
-\frac{1}{2\xi}(\nabla^\mu A^a_\mu)^2,
\]
together with the corresponding Faddeev--Popov ghost action. The resulting free field
operators remain normally hyperbolic, so the AQFT and pAQFT constructions of
Secs.~4--6 apply without modification.

\subsection{BRST symmetry}

The gauge--fixed theory admits a nilpotent BRST differential $s$ acting on the enlarged
field algebra. For gauge fields $A^a_\mu$, ghosts $c^a$, antighosts $\bar c^a$, auxiliary
fields $B^a$, and matter fields $\psi$, the BRST transformations take the standard form
\begin{align}
s A^a_\mu &= D_\mu c^a, \\
s c^a &= -\tfrac{1}{2} f^{abc} c^b c^c, \\
s \bar c^a &= B^a, \\
s B^a &= 0, \\
s \psi &= i g c^a T^a \psi,
\end{align}
with $s^2 = 0$ on all fields.

In the algebraic setting, physical observables are identified with BRST cohomology
classes. The interacting algebra constructed via pAQFT preserves BRST symmetry provided
the renormalization conditions respect the associated Ward identities.

\subsection{Renormalization and Ward identities}

The finite renormalization freedom in the definition of time--ordered products is
constrained by gauge symmetry. In particular, counterterms must be chosen so that the
quantum BRST operator remains nilpotent and the interacting theory satisfies the
Ward--Takahashi or Slavnov--Taylor identities.

Standard results ensure that this is always possible for renormalizable gauge theories
in four dimensions. Dimensional regularization with minimal subtraction provides one
explicit scheme in which these identities are preserved automatically, while the
covariant functional renormalization group realizes the same constraints at finite
coarse--graining scale.

\subsection{Unitarity and the physical Hilbert space}

Unitarity of the physical theory is recovered after restriction to the BRST cohomology.
Unphysical degrees of freedom introduced by gauge fixing (longitudinal modes, ghosts)
cancel from physical matrix elements.

In scattering regimes, unitarity of the $S$--matrix follows from the optical theorem,
which in turn is guaranteed by BRST invariance and the causal factorization property of
time--ordered products. In the algebraic formulation this unitarity statement is expressed
as the isometry of the $S$--matrix on the physical subspace of asymptotic states.

\subsection{Relation to Modal Triplet Theory}

From the Modal Triplet Theory perspective, gauge symmetry is not imposed but inherited
from internal isometries and bundle automorphisms of the modal geometry. The BRST
structure introduced here is therefore a technical device for quantization rather than a
fundamental ingredient. Its consistency and unitarity properties follow from the same
admissibility and stability conditions that guarantee the existence of the QFT
projection.

\subsection{Summary}

Gauge fixing, BRST symmetry, and unitarity are fully compatible with the AQFT and pAQFT
framework developed in this paper. They impose constraints on renormalization but do not
introduce additional assumptions or degrees of freedom. As a result, gauge theories
derived from Modal Triplet Theory admit a consistent and unitary perturbative quantum
description whenever the scattering assumptions of Section~8 are satisfied.
\section{Anomalies and global consistency}
\label{sec:anomalies}

Quantum consistency of gauge theories requires the absence of anomalies that would
invalidate gauge symmetry or diffeomorphism invariance at the quantum level. In this
section we show that the effective quantum field theory derived from Modal Triplet Theory
is anomaly--free, both locally and globally, and that this property follows from the
structure of the modal geometry rather than from ad hoc field--theoretic constraints.

\subsection{Local gauge anomalies}

Local (perturbative) gauge anomalies arise from the non--invariance of the quantum
effective action under infinitesimal gauge transformations. In four dimensions, they are
characterized by triangle diagrams involving chiral fermions.

In the present framework, the chiral fermion content of the effective theory is fixed by
the internal modal geometry and the ``two--of--three'' selection mechanism described in
Section~2. As a result, fermions appear in representations that satisfy the standard
anomaly cancellation conditions:
\begin{itemize}
\item pure gauge anomalies cancel within each gauge factor,
\item mixed gauge anomalies cancel between different factors,
\item mixed gauge--gravitational anomalies vanish.
\end{itemize}

These cancellations can be verified explicitly by computing the relevant traces over
representation matrices. However, their origin is geometric: anomaly--free combinations
correspond to admissible coherent sectors of the modal configuration space. No additional
anomaly cancellation conditions are imposed at the level of the effective field theory.

\subsection{BRST consistency and quantum gauge invariance}

Local anomaly cancellation ensures the existence of a nilpotent quantum BRST operator.
As discussed in Section~10, the renormalization freedom of perturbative algebraic quantum
field theory can be restricted so that the interacting theory satisfies the quantum BRST
identities (Slavnov--Taylor identities).

This guarantees that gauge invariance is preserved at the quantum level and that physical
observables, defined as BRST cohomology classes, are well defined. In particular, unitarity
of the physical $S$--matrix in scattering regimes follows from BRST invariance together
with causal factorization.

\subsection{Global anomalies}

In addition to local anomalies, certain gauge theories admit global (non--perturbative)
anomalies that are invisible in perturbation theory. The most prominent example in four
dimensions is the $\mathbb{Z}_2$ global anomaly of $SU(2)$ discovered by Witten.

In the present construction, the fermion content derived from Modal Triplet Theory always
contains an even number of left--handed $SU(2)$ doublets. This follows from the internal
bundle structure and family replication mechanism inherent in the modal geometry. As a
result, the mod--$2$ index governing the global $SU(2)$ anomaly vanishes, and the theory is
globally anomaly--free.

\subsection{Gravitational anomalies}

Pure gravitational anomalies do not occur in four dimensions for the field content under
consideration. Mixed gauge--gravitational anomalies are absent for the same reasons that
pure gauge anomalies cancel. The Hadamard state selected by the QFT projection ensures that
the renormalized stress--energy tensor is locally covariant and conserved, so that
diffeomorphism invariance is preserved at the quantum level.

\subsection{Structural origin of anomaly cancellation}

A key point of the present framework is that anomaly cancellation is not an independent
constraint imposed on the effective theory. Rather, it is a consequence of the existence
of a coherent, admissible modal sector admitting a stable QFT projection.

Modal configurations that would give rise to anomalous effective theories fail to satisfy
the admissibility and stability conditions required for coherent projection. In this sense,
anomaly cancellation is enforced ``upstream'' at the level of modal geometry rather than
``downstream'' at the level of quantum field theory.

\subsection{Summary}

The effective quantum field theory derived from Modal Triplet Theory is free of both local
and global anomalies. Gauge invariance, BRST consistency, and diffeomorphism invariance are
preserved at the quantum level without additional assumptions. This anomaly freedom is a
structural consequence of the modal geometry and the coherent sector selection, and it
ensures the internal consistency of the amplitudes and observables derived in earlier
sections.
\section{Confinement: what is and is not proven}
\label{sec:confinement}

A complete first--principles derivation of nonabelian confinement remains one of the
outstanding open problems in quantum field theory. In this section we clarify precisely
what can be established within the Modal Triplet Theory framework and the algebraic
quantum field theoretic construction developed in this paper, and what remains beyond its
current scope.

\subsection{Geometric origin of color degrees of freedom}

In Modal Triplet Theory, nonabelian gauge fields arise from internal symmetry bundles
associated with compact modal manifolds. In particular, color degrees of freedom are
supported on a distinguished internal bundle $B_3$ whose topology and holonomy structure
differ qualitatively from those associated with electroweak interactions.

The internal geometry of $B_3$ admits nontrivial cycles and flux sectors that do not
decompose into globally trivial configurations. As a result, color--charged excitations
cannot be represented as globally separable coherent modes on admissible slabs unless
they combine into color--singlet configurations.

This mechanism enforces color neutrality at the level of admissible coherent sectors and
is independent of perturbative dynamics.

\subsection{Wilson loops and area--law bounds}

Within the effective four--dimensional theory, the geometric structure of $B_3$ induces
nontrivial holonomy for nonabelian gauge fields. One can define Wilson loop observables
associated with closed spacetime curves $C$ and internal flux surfaces $S$.

Using standard inequalities, one obtains a lower bound of area--law type for the Wilson
loop expectation value,
\[
\langle W(C) \rangle \;\lesssim\; \exp[-\sigma\, \mathrm{Area}(S)] ,
\]
where the effective string tension $\sigma$ is determined by internal flux quantization
and geometric data of $B_3$. This bound reflects the energetic cost of separating
color--charged sources and provides a geometric obstruction to asymptotic color states.

The derivation of this bound relies only on geometric and topological features of the modal
bundle and does not depend on the details of perturbative renormalization.

\subsection{Limits of the present analysis}

Despite the existence of geometric area--law bounds and the absence of admissible
color--charged coherent sectors, the present framework does \emph{not} constitute a full
nonperturbative proof of confinement in the sense of Yang--Mills theory.

In particular:
\begin{itemize}
\item a complete control of the nonabelian Wilson loop at all scales is not provided,
\item the existence of a mass gap in pure Yang--Mills theory is not proven here,
\item the detailed infrared dynamics of flux tube formation and breaking are not derived.
\end{itemize}

These limitations are shared by all perturbative and semiclassical approaches to
confinement and are not specific to Modal Triplet Theory.

\subsection{Interpretation within the amplitudes framework}

From the perspective of the present amplitudes construction, confinement enters as a
\emph{consistency constraint} rather than as a dynamical theorem. Scattering amplitudes
involving colored external states are excluded because such states do not correspond to
admissible asymptotic particle configurations. Only color--singlet combinations admit
asymptotic states and well--defined scattering amplitudes.

This restriction is compatible with all perturbative calculations presented earlier:
whenever perturbative QCD amplitudes are computed, they are implicitly understood as
describing short--distance processes embedded in color--neutral hadronic states.

\subsection{Summary}

The Modal Triplet Theory framework provides a geometric mechanism that enforces color
neutrality and yields area--law bounds suggestive of confinement. While this falls short of
a complete nonperturbative proof of Yang--Mills confinement, it explains why colored
asymptotic states are excluded and why scattering amplitudes must be formulated in terms
of color--singlet degrees of freedom.

In this sense, confinement is built into the admissibility structure of the theory rather
than emerging as a purely infrared dynamical phenomenon.
\section{Synthesis with the Modal Triplet Theory corpus}
\label{sec:synthesis}

We now place the results of this paper within the broader Modal Triplet Theory (MTT)
program. The purpose of this section is not to introduce new technical content, but to
clarify how the amplitudes construction developed here interfaces with, and completes,
the existing MTT corpus.

\subsection{Position of the amplitudes layer}

The MTT program may be viewed as consisting of a sequence of logically distinct layers:
\begin{enumerate}
\item a foundational layer establishing coherent fixed points, spectral gaps, and
admissible domains,
\item a kinematic layer defining localization, propagation, and persistence without
assuming background spacetime notions,
\item a quantum field theory layer deriving algebraic QFT on curved spacetimes,
\item an amplitudes layer deriving perturbative scattering amplitudes and phenomenology.
\end{enumerate}

The present paper occupies the final position in this sequence. Its role is to show how,
once an algebraic quantum field theory has been obtained from Modal Triplet Theory, the
familiar perturbative machinery of quantum field theory emerges in a mathematically
controlled and regime--aware manner.

\subsection{Relation to AQFT nets from admissible charts}

In companion work, algebraic quantum field theory nets are derived directly from the
overlap structure of admissible charts in Modal Triplet Theory, without assuming a
background spacetime or particle interpretation. When the QFT projection exists, those
admissibility--indexed nets reduce to the spacetime--indexed Haag--Kastler nets used in
the present paper.

Sections~4 and~5 of this work can therefore be read as a specialization of the more general
net construction to regimes admitting an effective spacetime description. The amplitudes
derived here are consequently conditional on the same admissibility criteria that govern
the existence of the AQFT net itself.

\subsection{Relation to QFT on curved spacetime}

The derivation of algebraic quantum field theory on curved spacetime from Modal Triplet
Theory provides the immediate input for the present work. In that construction, the
projection $\Pi_{\mathrm{QFT}}$ yields:
\begin{itemize}
\item a locally covariant field algebra,
\item a Hadamard state ensuring renormalized local observables,
\item a controlled description of scale dependence via the covariant functional
renormalization group.
\end{itemize}

The present paper does not modify or extend that construction. Instead, it addresses the
question of how perturbative amplitudes and standard phenomenology arise from it when
additional scattering assumptions are satisfied. In this sense, the amplitudes layer is
strictly downstream of the curved--spacetime QFT layer.

\subsection{Relation to kinematics and irreversibility}

Earlier work within the MTT corpus has emphasized that notions of position, motion, and
irreversibility are not fundamental, but arise from the persistence of coherent structures
across overlapping admissible descriptions. The algebraic propagation and time--slice
axiom used in the present work are consistent with this perspective: propagation is defined
without reference to trajectories, and irreversibility arises from the absence of global
sections rather than from microscopic nonunitarity.

Scattering amplitudes therefore represent a further specialization of kinematic
persistence, valid only when coherent structures admit asymptotic free descriptions. This
interpretation aligns the amplitude concept with the general MTT view that physical notions
are conditional rather than absolute.

\subsection{Completeness of the program}

Taken together with the foundational, kinematic, and QFT projection papers, the present
work completes the Modal Triplet Theory program at the level of perturbative quantum field
theory. All standard structures---locality, propagation, interactions, renormalization,
and amplitudes---are derived from modal geometry and coherence, with their regimes of
validity made explicit.

No additional postulates are required beyond those already present in the MTT foundation.
Conversely, no claim is made that perturbative amplitudes exhaust the physical content of
the theory. Nonperturbative phenomena, global questions, and strong--coupling dynamics
remain subjects for future investigation.

\subsection{Summary}

The amplitudes construction developed in this paper should be understood as the final,
downstream realization of Modal Triplet Theory in regimes where perturbative quantum field
theory is applicable. It neither replaces nor competes with the earlier layers of the
program, but rather shows how they culminate in the familiar computational framework used
in particle physics when appropriate conditions are met.
\section{Conclusions and outlook}
\label{sec:conclusions}

In this paper we have completed the derivation of perturbative quantum field theory
amplitudes from Modal Triplet Theory (MTT) in a mathematically rigorous and logically
closed manner. Starting from the coherent fixed--point sector of the modal configuration
space, we have shown how the full perturbative machinery of quantum field theory emerges
as a downstream consequence of algebraic and geometric structures, rather than as a set
of independent postulates.

The central results of this work may be summarized as follows.

First, given an admissible coherent modal sector, the MTT $\rightarrow$ QFT projection
yields a locally covariant algebraic quantum field theory on a globally hyperbolic
four--dimensional spacetime, equipped with a Hadamard state. This provides a well--defined
framework for local observables, causal propagation, and renormalization on curved
backgrounds.

Second, interactions are constructed from first principles using perturbative algebraic
quantum field theory. Interacting fields are defined via time--ordered products and the
Bogoliubov map, with renormalization appearing as a finite, local ambiguity constrained by
covariance, power counting, and gauge symmetry. Standard schemes such as dimensional
regularization with minimal subtraction and the covariant functional renormalization group
are recovered as concrete realizations of this algebraic structure.

Third, the familiar Feynman rules of perturbative quantum field theory are shown to arise
from the Wick expansion of time--ordered products in Hadamard states. Diagrams, propagators,
and vertices are therefore derived artifacts of the algebraic construction, not primitive
assumptions.

Fourth, scattering amplitudes are defined precisely and conditionally. Using Haag--Ruelle
and LSZ theory, amplitudes emerge only in regimes admitting asymptotic particle states,
such as asymptotically flat or stationary spacetimes. In generic curved or time--dependent
backgrounds, the same formalism yields local or in--in observables rather than a global
$S$--matrix. This distinction is essential for a correct interpretation of quantum field
theory in cosmological and gravitational settings.

Fifth, all effective couplings, masses, mixing parameters, and thresholds entering the
amplitudes are determined by bounded overlap integrals and curvature--gap data intrinsic
to Modal Triplet Theory. When evaluated with identical low--energy inputs, the resulting
amplitudes coincide with those of the Standard Model, while retaining a clear geometric
origin in the modal structure.

Taken together with earlier work on the foundations, kinematics, and curved--spacetime
quantum field theory projection of MTT, the present paper completes the amplitudes layer
of the program. Locality, propagation, interactions, renormalization, and perturbative
phenomenology are all shown to arise from modal geometry and coherence, with their domains
of validity stated explicitly.

Several important directions remain open. A full nonperturbative treatment of confinement
and strong--coupling dynamics lies beyond the scope of the present work. Global--in--time
questions, including cosmological initial conditions and singularity resolution, require
further development of the admissibility and fixed--point framework. Extensions to finite
temperature, nonequilibrium steady states, and genuinely quantum gravitational regimes
remain to be explored.

These open problems do not detract from the main conclusion of this paper. Rather, they
underscore the central lesson of Modal Triplet Theory: physical structures are conditional
on coherence and admissibility, and familiar concepts such as particles, amplitudes, and
scattering emerge only in regimes where those conditions are satisfied.

The derivation presented here demonstrates that perturbative quantum field theory is not
a fundamental starting point, but a consequence of deeper geometric and algebraic
principles. In this sense, Modal Triplet Theory provides a unified and conceptually
transparent path from modal geometry to the full perturbative apparatus of modern
particle physics.
\appendix
\section{Canonical normalization and overlap integrals}

This appendix collects the technical details underlying the canonical normalization of
the effective four--dimensional fields and the evaluation of overlap integrals used to
define couplings in Section~3.

Let $\omega_a$ denote a harmonic representative on the internal bundle $B_a$, normalized
with respect to the $L^2$ inner product induced by the internal metric:
\[
\int_{B_a} \langle \omega_a,\omega_a\rangle\, d\mu_{B_a} = 1 .
\]
With this choice, dimensional reduction of the higher--dimensional kinetic terms yields
canonically normalized gauge, fermion, and scalar kinetic terms on $(Y_4,g)$.

All numerical prefactors arising from the higher--dimensional action and internal volumes
are absorbed into fixed constants $N_a$ multiplying the overlap integrals. No further
field rescalings are required in the effective theory.
\section{Propagator and vertex rule cards}

For convenience, we collect here the explicit propagators and vertices corresponding to
the effective action of Section~3. These expressions are not postulated, but arise from
the graph expansion of time--ordered products derived in Section~7.

Metric signature is $(+,-,-,-)$ and all momenta are taken incoming.

\paragraph{Propagators (Feynman gauge).}
\begin{itemize}
\item Scalar: $\displaystyle \frac{i}{p^2 - m^2 + i0}$,
\item Fermion: $\displaystyle \frac{i(\slashed{p}+m)}{p^2 - m^2 + i0}$,
\item Gauge boson: $\displaystyle \frac{-ig_{\mu\nu}}{p^2 + i0}$.
\end{itemize}

\paragraph{Vertices.}
\begin{itemize}
\item QED: $-ie\gamma^\mu$,
\item QCD: $-ig_s\gamma^\mu T^a$,
\item Yukawa: $-iY_{ij}$,
\item Three--gauge vertex: standard Yang--Mills form.
\end{itemize}

These rules coincide with the Standard Model when evaluated with identical low--energy
inputs.
\section{Representative decay calculations}

This appendix provides explicit formulas used in the benchmark calculations of
Section~9.

\paragraph{Muon decay.}
\[
\Gamma_\mu =
\frac{G_F^2 m_\mu^5}{192\pi^3}
f\!\left(\frac{m_e^2}{m_\mu^2}\right)
\left[1+\Delta_{\mathrm{RC}}^{(\mu)}\right],
\quad
f(x)=1-8x+8x^3-x^4-12x^2\ln x .
\]

\paragraph{Pseudoscalar leptonic decays.}
\[
\Gamma(P\to \ell\nu[\gamma]) =
\frac{G_F^2 |V_{qQ}|^2 f_P^2 m_\ell^2 m_P}{8\pi}
\left(1-\frac{m_\ell^2}{m_P^2}\right)^2
\left[1+\delta_{\mathrm{EM}}+\delta_{\mathrm{SU(2)}}\right].
\]

All quantities are evaluated using overlap--determined couplings or empirical inputs, as
described in the main text.
\section{Gauge coupling normalization from modal geometry}

Gauge couplings are defined by quadratic overlap integrals on internal bundles. For
hypercharge we adopt the GUT--normalized convention
\[
g_1 = \sqrt{\frac{5}{3}}\, g' ,
\]
so that the one--loop coefficient takes its standard value.

A representative example on $S^1_{\mathrm{cen}}$ yields
\[
\|\omega\|_{L^2}^2 = \int_{0}^{2\pi R} R^{-2} R\, d\theta = 2\pi ,
\]
fixing the normalization uniquely.
\section{BRST symmetry and Ward identities}

Gauge fixing introduces ghost fields and a nilpotent BRST differential $s$. Physical
observables are identified with BRST cohomology classes.

Renormalized time--ordered products can always be chosen so that the interacting theory
satisfies the quantum BRST identities. In dimensional regularization with minimal
subtraction, Ward identities are preserved automatically.

This guarantees unitarity of the physical $S$--matrix in scattering regimes.
\section{Curvature--gap coefficients and threshold matching}

Representation--dependent curvature corrections take the form
\[
M^2 = m^2 + \beta_{\mathrm{rep}}\, \mathrm{Scal} + \cdots ,
\]
with
\[
\beta_{\mathrm{spinor}}=\frac{1}{4},
\qquad
\beta_{\mathrm{scalar}}=\frac{1}{6}.
\]

These coefficients determine physical thresholds entering renormalization group flow
and decoupling.
\section{Wilson loops and geometric confinement bounds}

Let $W(C)$ denote a Wilson loop associated with a closed curve $C$. For nontrivial internal
flux sectors one obtains
\[
\langle W(C) \rangle \lesssim \exp[-\sigma\, \mathrm{Area}(S)],
\]
where $S$ is a spanning surface and $\sigma$ depends on internal geometric data.

This provides a lower bound consistent with confinement, but does not constitute a full
nonperturbative proof.

\bigskip
%\nocite{*}
\bibliographystyle{plain} % or abbrv, unsrt, alpha, etc
\bibliography{main}

\end{document}
\typeout{get arXiv to do 4 passes: Label(s) may have changed. Rerun}